\documentclass[10pt,a4paper]{article}

\usepackage[margin=1in]{geometry}
\usepackage{amsmath,amssymb,amsthm}
\usepackage{mathtools}
\usepackage{booktabs}
\usepackage{array}
\usepackage{longtable}
\usepackage{hyperref}
\usepackage{xcolor}

\hypersetup{colorlinks=true, linkcolor=blue!50!black, citecolor=blue!50!black, urlcolor=blue!50!black}

\title{Negative Results for Elliptic-Curve Index Calculus:\\
       Six Statements, the Proofs or Evidence Behind Them,\\
       and What They Do Not Say}
\author{Anonymous}
\date{}

\newcommand{\FF}{\mathbb{F}}
\newcommand{\ZZ}{\mathbb{Z}}
\newcommand{\Aut}{\mathrm{Aut}}
\newcommand{\ev}[1]{\textup{[#1]}}

\newtheorem{theorem}{Theorem}
\newtheorem{proposition}[theorem]{Proposition}
\newtheorem{lemma}[theorem]{Lemma}
\newtheorem{corollary}[theorem]{Corollary}
\theoremstyle{definition}
\newtheorem{definition}[theorem]{Definition}
\theoremstyle{remark}
\newtheorem{remark}[theorem]{Remark}

\begin{document}
\maketitle
\emergencystretch=2em

\begin{abstract}
We collect six negative results on index calculus for the elliptic-curve
discrete logarithm and state each at the tier its evidence supports.
(1) The compact-orbit Koblitz decomposition configurations measured in this
programme sit on or above the generic preprocessing line
$ST^2 \gtrsim r/|\mathrm{Aut}|$, and at the rank-stage mean on the
non-adaptive line $ST \approx r/|\mathrm{Aut}|$.
(2) Isogeny oracles add nothing in a generic group model: we reprove the bound
$(Q+2)^2/r + 1/r$ and give a Lean statement of its counting core.
(3) The mean decomposition yield of a factor base of size $B$ is
$\binom{B+m-1}{m}/N$ whatever its geometry.
(4) The monodromy group of the Semaev summation cover is
$(\mathbb{Z}/2)^{m-2}$; this is a derivation that has not had an independent
review.
(5) Lemma 4 of Nagao's ePrint 2015/984, $d_F \le d'_F$, is false over
$\mathbb{F}_2$ under the equality reading of its definitions; one witness is
checked by hand and both are re-derived natively, and nothing is claimed about the
paper's Proposition 5.
(6) We could not establish that Semaev's Assumption 1 fails off the diagonal,
and state what is and is not known. Every number links to a committed
evidence file. None of this is a speedup or a statement about a deployed curve.
\end{abstract}

\section{Introduction}
\label{sec:intro}

A programme that searches for faster index calculus on elliptic curves
produces many negative results, and they are easy to lose: each sits in a
different repository, at a different strength, behind a different convention.
This paper states six of them in one place. For each we say what is proved, what
is an exact finite computation, what is a measured committed record, and what is
only a derivation that nobody outside the programme has read. We also say what
each statement does \emph{not} say.

The six statements are:
\begin{description}
\item[S1] compact-orbit Koblitz index calculus is not evidence of a non-generic
  speedup: its measured configurations sit on or above the generic
  preprocessing lines (Section~\ref{sec:s1});
\item[S2] isogeny oracles add nothing in the generic group model
  (Section~\ref{sec:s2});
\item[S3] the mean decomposition yield is conserved under factor-base geometry
  (Section~\ref{sec:s3});
\item[S4] the monodromy of the Semaev summation cover is $(\ZZ/2)^{m-2}$
  (Section~\ref{sec:s4});
\item[S5] Lemma 4 of Nagao's ePrint 2015/984 is false as a universal statement
  (Section~\ref{sec:s5});
\item[S6] whether Semaev's 2015 Assumption 1 fails off the diagonal is
  \emph{not} established by the programme's records (Section~\ref{sec:s6}).
\end{description}
S6 is included because the design of this paper listed it as a negative result,
and the evidence on disk does not support that. Reporting it as open is the
result.

\paragraph{Claim class.}
Under the programme's success criterion a relabelling or an accounting change is
not an improvement, and this paper is that: a consolidation of existing
negatives. It measures no new speed, reports no gap to Pollard rho of its own
(the studies it cites report theirs), and makes no key-recovery, sub-rho or
deployed-curve claim.

\paragraph{Labels.}
Each statement carries one label, fixed by its evidence and not by our
judgement of how convincing it is.
\textsc{proved}: a proof is in the paper.
\textsc{exact}: an exact finite computation that we re-ran natively, with its
scope stated.
\textsc{recorded}: committed measured rows, cited by file and digest.
\textsc{derived, unreviewed}: a written derivation that has not had an
independent review.
\textsc{not established}: the records do not support the statement.
Table~\ref{tab:summary} lists them.

\begin{table}[t]
\centering\small
\begin{tabular}{@{}l p{4.6cm} p{2.9cm} p{4.4cm}@{}}
\toprule
 & statement & label & evidence \\
\midrule
S1 & compact-orbit IC on or above $ST^2 \gtrsim r/|\Aut|$ & \textsc{recorded} (14 rows) & \ev{E1}, \ev{E2}, \ev{E3} \\
S2 & isogeny oracles add nothing generic & \textsc{proved}; counting core in Lean (not yet kernel-checked) & \ev{E4}, \ev{E31} \\
S3 & mean yield conserved under base geometry & \textsc{proved}; \textsc{exact}; \textsc{recorded} & \ev{E5}, \ev{E7}, \ev{E6}, \ev{E18}, \ev{E32} \\
S4 & Semaev cover monodromy $(\ZZ/2)^{m-2}$ & \textsc{derived, unreviewed} & \ev{E9}, \ev{E10}, \ev{E11}, \ev{E13} \\
S5 & Nagao Lemma 4 false (equality reading, $\FF_2$) & \textsc{exact} & \ev{E14}, \ev{E15}, \ev{E18} \\
S6 & Assumption 1 fails off the diagonal & \textsc{not established} & \ev{E20}, \ev{E21}, \ev{E22} \\
\bottomrule
\end{tabular}
\caption{The six statements, their labels and evidence. Identifiers $E_i$ are
rows of the evidence ledger (Appendix~\ref{app:ledger}).}
\label{tab:summary}
\end{table}

\paragraph{Conventions.}
Evidence files live in two repositories and are cited with the prefix
\texttt{crypto:} or \texttt{crypto-autoresearcher:};
Appendix~\ref{app:ledger} gives each file's SHA-256. $r$ is the prime order of
the subgroup under attack, $A=|\Aut|$ is the order of the automorphism group the
base is folded by, $S$ is the stored advice (table entries, not bits), $T$ is the
online cost in the unit named in each section. All measurements are at toy
scale. A \emph{probe} (one quadratic solve plus a root-table lookup) is not a
group operation; where a bound is stated in group operations the section says
so.

\section{S1: compact-orbit Koblitz index calculus against the generic lines}
\label{sec:s1}

\paragraph{Setting.}
The compact-orbit configuration on a Koblitz curve over $\FF_{2^n}$ folds a
factor base by the Frobenius orbit, collects $4$-sum relations through a table of
roots of the third summation polynomial, and extracts a target by a fixed
state-and-shift scan of that table. The table has $S=K^2 n$ entries. The study
we rely on (\ev{E3}, \ev{E1}, \ev{E2}) audits the committed records of this
configuration for $n\in\{37,41,53,61,71,73,83\}$ against three published lower
bounds, and we use only its \emph{adaptive} and \emph{non-adaptive}
preprocessing lines:
\begin{itemize}
\item the adaptive line $S\,T^2 \approx r$ (Bernstein--Lange; Corrigan-Gibbs and
  Kogan, which also gives the matching algorithm) \cite{BL13,CGK18};
\item the non-adaptive line: a non-adaptive algorithm with advice $S$ and $T$
  queries succeeds with probability at most $3T^2/r + 4\ln 2\cdot ST/r$, so
  below $T=\sqrt r$ it needs $ST \gtrsim r$ \cite{DKM25}.
\end{itemize}
Dividing $r$ by $A$ is the audit's own convention (\ev{E3}: ``Theorem 3.2 has
none''), not a theorem of the cited papers. We keep that division and say so.

\paragraph{Result.}
Table~\ref{tab:s1} reproduces the audit's compact-orbit ladder. The first block
uses the single frozen target probes; the second the rank-stage mean probes per
relation.

\begin{table}[t]
\centering\small
\begin{tabular}{@{}lrrrrrr@{}}
\toprule
rung & $\log_2 r$ & $A$ & $S=K^2n$ & $\log_2(ST^2)-\log_2(r/A)$ & $ST/(r/A)$ & $ST/(r/A)$ \\
 & & & & frozen target & frozen target & rank-stage mean \\
\midrule
n37 & 27.78 & 74  & $1.813\times10^{3}$ & 9.4  & 0.627 & 1.094 \\
n41 & 39.00 & 82  & $2.962\times10^{5}$ & 14.3 & 0.939 & 1.092 \\
n53 & 44.26 & 106 & $2.565\times10^{6}$ & 13.1 & 0.342 & 1.081 \\
n61 & 47.21 & 122 & $2.196\times10^{7}$ & 19.7 & 3.741 & 1.053 \\
n71 & 52.29 & 142 & $2.556\times10^{7}$ & 7.1  & 0.010 & 1.478 \\
n73 & 56.26 & 146 & $2.628\times10^{7}$ & 13.2 & 0.020 & 2.734 \\
n83 & 52.93 & 166 & $2.988\times10^{7}$ & 25.4 & 5.120 & 1.417 \\
\bottomrule
\end{tabular}
\caption{Compact-orbit ladder (\ev{E2}, rows \texttt{target\_probes\_online} and
the rank-stage-mean sentence; same numbers in \ev{E1}). $T$ is in probes, not
group operations.}
\label{tab:s1}
\end{table}

\begin{proposition}[S1, \textsc{recorded}]
\label{prop:s1}
On the seven rungs of Table~\ref{tab:s1}:
(a) every measured row, whether the single frozen target or the rank stage plus
target, has $\log_2(ST^2)-\log_2(r/A)\ge 7.1$ (frozen target: $7.1$ to $25.4$;
rank stage plus target: $16.8$ to $45.8$);
(b) at the rank-stage mean probes per relation, $ST/(r/A)$ lies in
$[1.053,\,2.734]$;
(c) for single frozen targets $ST/(r/A)<1$ at n37, n41, n53, n71 and n73.
\end{proposition}

The audit also records, for the whole study, that none of $208$ configurations
(prime, random binary, Koblitz round five, one-target and batch Koblitz, and the
compact-orbit probe rows) has a measured cost below its bound (\ev{E3}, Result 1;
smallest measured-to-bound ratios $38.9$, $56.8$, $355$, $2.84\times10^{6}$,
$7.99\times10^{4}$ and $1077$ for the six families).

\paragraph{How to read (c).}
The non-adaptive bound limits the \emph{success probability} of an algorithm over
random targets; a single cheap target is not a counterexample. The audit finds the
two cheapest draws, n71 and n73, at $0.0065$ and $0.0075$ times the rank-stage mean
probes per relation; under an exponential first-hit law each is a
$P\approx0.006$--$0.007$ event, and it calls them improbably cheap (\ev{E3},
``Record problems'', item 3). We report them and do not average them away.

\paragraph{What S1 does not say.}
It is not a proof that the configuration is a generic algorithm. The structured
generic-group bound is vacuous here: every rung has $\delta_{\mathrm{upper}}=1$
(its $4$-sums cover the subgroup), so it certifies nothing (\ev{E3}, Result 1).
The comparison is in probes and table entries, and the audit flags the rows
\texttt{unit\_not\_group\_operations}. The audit also found three record problems
that we carry forward: the n61 subgroup order is recorded as $\log_2 r=47.54$ in
the ledger against the exact $47.2109$; the excluded-precompute field means
different things at different rungs; and the n71 and n73 frozen targets are
improbable (\ev{E3}, ``Record problems''). The distance to matched rho at full
cost is outside this paper.

\section{S2: isogeny oracles add nothing in the generic group model}
\label{sec:s2}

Let $E_1,\dots,E_k$ be curves in one $\FF_q$-isogeny class with $N=h\,r$, $r$
prime and $r\nmid h$. Each order-$r$ subgroup $G_j\subseteq E_j(\FF_q)$ is encoded
by an independent random injection $\sigma_j:\ZZ/r\to\{0,1\}^{*}$. The algorithm
has oracles for addition and negation in every $G_j$ and for every
$\FF_q$-rational isogeny $\psi:E_i\to E_j$ of any degree, restricted to $G_i$,
and is given the scalars $\lambda_\psi\in\ZZ/r$ through which each $\psi$ acts.

\begin{theorem}[Theorem 15.1 of \ev{E4}; \textsc{proved}]
\label{thm:151}
An algorithm that makes $Q$ oracle calls finds $\log_{g_1}(Q_1)$ with
probability at most $(Q+2)^2/r + 1/r$.
\end{theorem}

\paragraph{Counting convention.}
The note does not say how the generators of the other $G_j$ are supplied. We read
the theorem as giving the algorithm the encodings of $g_1$ and $Q_1$ only, with
every other element obtained through an oracle answer; then at most $Q+2$ distinct
elements ever exist. If encodings of fixed generators of the other groups are
also given, $Q+2$ becomes $Q+k+1$ below and nothing else changes.

\paragraph{Proof.}
Write $x=\log_{g_1}Q_1$, uniform in $\ZZ/r$. Because $G_i$ has prime order,
$\ker\psi\cap G_i$ is $0$ or $G_i$, so a rational isogeny maps $G_i$ either
isomorphically onto $G_j$ or to the identity (this includes isogenies of degree
divisible by $r$ whose kernel is the other Frobenius eigenline of $E[r]$). With
the $\lambda_\psi$ known, every element the algorithm can hold has the form
$\sigma_j(a+bx)$ with $(a,b)\in(\ZZ/r)^2$ computed from the query structure and
the $\lambda_\psi$ alone: group operations add and negate forms, and an isogeny
multiplies a form by $\lambda_\psi$ or by $0$.

The argument is Shoup's \cite{Sho97}, adapted to forms on several curves.
Game~0 is the real game. Game~1 is run by a simulator that tracks only the forms
$(a,b)$, giving the forms $(1,0)$ and $(0,1)$ to the inputs $g_1$ and $Q_1$. It answers each query with the string already assigned to the same form
on the same curve, and with a fresh random string if the form is new. The simulator
never uses $x$. The two games give the same transcript unless two \emph{distinct}
forms on one curve have the same value at $x$, the bad event. In Game~1 the
algorithm's view is independent of $x$, so it outputs the right $x$ with
probability exactly $1/r$.

It remains to bound the bad event. Fix the simulator's random choices; they fix the
families $P_1,\dots,P_k$ of forms issued on each curve, with
$\sum_j |P_j|\le Q+2$. The set of secrets at which two different members of $P_j$
coincide has at most $|P_j|^2-|P_j|$ elements, by Lemma~\ref{lem:core} below, which
holds over any field. Summing over $j$ gives at most $(Q+2)^2$ bad secrets out of
$r$. Averaging over the simulator's choices, the bad event has probability at most
$(Q+2)^2/r$, so the real algorithm wins with probability at most
$1/r+(Q+2)^2/r$. $\square$

\begin{lemma}[counting core]
\label{lem:core}
Let $F$ be a field and $P\subseteq F\times F$ a finite set of affine forms
$(a,b)$, read as $x\mapsto a+bx$. (i) Two distinct forms agree at no more than one
$x\in F$. (ii) For any finite $S\subseteq F$, the number of $x\in S$ at which two
different members of $P$ agree is at most $|P|^2-|P|$. (iii) Off that set the
evaluation map $p\mapsto a+bx$ is injective on $P$.
\end{lemma}
\begin{proof}
(i) If $a+bx=c+dx$ and $a+by=c+dy$ with $(a,b)\neq(c,d)$: if $b=d$ then $a=c$, a
contradiction; otherwise $(b-d)(x-y)=0$ gives $x=y$. (ii) The bad set lies in the
union, over the $|P|^2-|P|$ ordered pairs of distinct forms, of the sets in (i).
(iii) is the definition. \qedhere
\end{proof}

The bound $|P|^2-|P|$ counts ordered pairs; the unordered count $\binom{|P|}{2}$
holds as well, and is attained (see below).

\paragraph{Machine support.}
Lemma~\ref{lem:core} is stated and proved in Lean in
\texttt{IcNegative/AffineCore.lean} (\ev{E31}; Section~\ref{sec:lean}), and it is
checked exhaustively on $\ZZ/p$ for $p=5,7,11$ in the native instrument
(\ev{E18}): $2352$ ordered pairs of distinct forms on $\ZZ/7$ and $14520$ on
$\ZZ/11$ agree at most once; over all $18424$ sets of three forms on $\ZZ/7$ the
largest bad set has $3=\binom32$ elements, and over all $211876$ sets of four forms
on $\ZZ/7$ it has $6=\binom42$. The negative control behaves as it must: on $\ZZ/8$
the forms $(0,0)$ and $(0,4)$ agree at $4$ points, and on $\ZZ/9$ the forms
$(0,0)$ and $(0,3)$ at $3$, so the field hypothesis is needed.

\paragraph{What S2 does not say.}
Theorem~\ref{thm:151} concerns isogeny \emph{maps} used as oracles. Two things sit
outside its hypotheses on purpose (\ev{E4}): the Weil and Tate pairing on the
$[r]$-fibre is not an isogeny oracle, and Smart's attack uses a lift to $\mathbb{Q}_p$.
Any effect of an isogeny on index calculus must therefore come from the
\emph{representation}, the coordinates in which factor bases and summation
polynomials live. The Lean statement is the counting step only; the oracle model,
the simulator and the hybrid step are not formalised (\texttt{OPEN\_TARGETS.md}, O1).

\section{S3: mean decomposition yield is conserved under base geometry}
\label{sec:s3}

Let $G$ be a finite abelian group of order $N$ and $D\subseteq G\setminus\{0\}$ a
factor base of $B$ distinct elements. For $m\ge1$ and $g\in G$ let $c_D(g)$ be the
number of size-$m$ multisets from $D$ whose sum is $g$.

\begin{proposition}[yield conservation, \textsc{proved}]
\label{prop:yield}
$\sum_{g\in G} c_D(g)=\binom{B+m-1}{m}$. Hence the mean per-target yield is
$\binom{B+m-1}{m}/N$ for every base of size $B$, and the fraction of targets with
at least one decomposition is at most $\min\{1,\ \binom{B+m-1}{m}/N\}$.
\end{proposition}
\begin{proof}
Every size-$m$ multiset has exactly one sum, so the multisets are partitioned by
their sums; there are $\binom{B+m-1}{m}$ of them. The coverage bound is the
pigeonhole bound on the image of a map from that set to $G$. \qedhere
\end{proof}

The identity holds for any map from a finite domain, not only for multiset sums.
Lean (\ev{E32}): \texttt{sum\_yield} is $\sum_g \mathrm{yield}(\varphi,g)=|\tau|$,
\texttt{total\_yield\_eq} says two maps on domains of equal size have equal total
yield, and \texttt{coverage\_le\_card\_dom} and \texttt{coverage\_le\_card\_cod}
bound the number of covered targets by $\min(|\tau|,|G|)$. The list form of the same
identity (\texttt{sum\_count}, \texttt{patterns\_neutral}) is already kernel-checked
in \ev{E7}, with the axiom log \ev{E8} showing only the standard foundations. The
binomial count of the domain $\mathrm{Sym}^m(D)$ is not formalised (O2).

\paragraph{Corollary (screening rule).}
A proposal that promises a higher \emph{mean} yield at matched base size is refuted
by arithmetic before any experiment. Only coverage, multiplicity, relation rank,
recognisability and the cost of finding a decomposition are free.

\paragraph{Exact and recorded support.}
(i) The programme's finding records that over $144$ cells at $N\sim2^{14},2^{16},2^{18}$
on $12$ generated prime-order curves with $m=3$ the largest absolute deviation of
the exact cell mean from $\binom{B+2}{3}/N$ was exactly $0$ (\ev{E5}, Consequence 1).
(ii) The native instrument enumerates every base of size $B\in\{3,4\}$ in
$\ZZ/N$ for $N\in\{11,12,13\}$ ($1540$ bases in all) and finds the total always equal
to $\binom{B+2}{3}$, while coverage varies, for example between $3$ and $9$ targets
of $12$ at $B=3$, $N=12$, and an arithmetic-progression base has the same total
with coverage below the maximum (\ev{E18}). (iii) The cyclotomic $\gamma$-weight base
was tested against $1999$ random Frobenius-stable bases of equal column count at
$n=23$: on all $10$ coverage tests and all $5$ energy tests it lies inside the null
distribution, smallest two-sided $p=0.363$ against a critical $0.00333$, largest
deviation from the null median $0.65\%$, while the positive control is detected
at the floor $p=0.001$ (\ev{E6}). (iv) At toy sizes of $20$ to $30$ bits and $m=2$, the matched-size yield ratios of
the $y$-window arms against a random set of the same size lie inside $[0.95,1.03]$
(\ev{E33}).

\paragraph{What S3 does not say.}
It does not say that structured bases cannot beat random bases: coverage headroom
up to $\min\{1,\mu\}$ is real (a Bose--Chowla Sidon base reaches it, \ev{E5}), and
under a one-relation-per-target solve the gain is at most
$\min\{1,\mu\}/(1-e^{-\mu})\le1.582$. It says nothing about the cost of finding a
decomposition or the linear-algebra stage, where the cost of index calculus sits.
The yield-neutral result of \ev{E6} tests the yield clause only; the solving-degree
clause of that lead was not tested there. In the campaign of 2026-10-18 the planned
toy check of conservation inside the kernel-program search (item 10) has not run
(\ev{E23}); we cite no number from it. Item 17 (\ev{E35}) measures the counted cost of a cyclotomic base at $n=83$, not its yield: it reports that the base-dependent part of the cost is small next to the arm-independent pair table, and it is context for S3, not evidence for or against it.

\section{S4: the monodromy of the Semaev cover is $(\ZZ/2)^{m-2}$}
\label{sec:s4}

Let $E$ be an elliptic curve over any field $K$, $\pi:E\to\mathbb{P}^1$ the
quotient by inversion, $m\ge3$ and $n=m-1$. Let $t_1,\dots,t_n$ be independent
transcendentals, $L=K(t_1,\dots,t_n)$, and $P_i$ the tautological point with
$\pi(P_i)=t_i$. The $m$-th Semaev summation polynomial $S_m(t_1,\dots,t_n,T)\in L[T]$
has degree $2^{m-2}$ in $T$, and its roots are the $x$-coordinates of the signed
sums $\epsilon_1P_1+\dots+\epsilon_nP_n$.

\begin{theorem}[\ev{E9}, \ev{E10}; \textsc{derived, unreviewed}]
\label{thm:mono}
\emph{(A)} $S_m$ is separable, its splitting field is the function field of
$E^{n}/\Delta$ with $\Delta$ the diagonal inversion, and
$\mathrm{Gal}(S_m/L)=(\ZZ/2)^{n}/\Delta\cong(\ZZ/2)^{m-2}$, acting simply
transitively on the $2^{m-2}$ roots; arithmetic and geometric monodromy agree. The
group is the same for every curve, base field and characteristic.
\emph{(B)} Over $\FF_q$, at every specialisation outside an explicit degenerate
locus of size $O_m(q^{m-2})$, $S_m(a,T)$ is squarefree and is either a product of
$2^{m-2}$ distinct linear factors, when the $n$ fibres of $\pi$ over the $a_i$ are
simultaneously rational, or a product of $2^{m-3}$ irreducible quadratics; no
factor of degree $\ge3$ and no mixed type occurs.
\emph{(C)} If every $a_i$ is the $x$-coordinate of an $\FF_q$-rational point, $S_m(a,T)$
splits completely, at every $m$: the factor-base locus lies inside the totally
split locus.
\end{theorem}

\paragraph{Derivation.}
The roots are the $x$-coordinates of the signed sums. The splitting field sits
inside $K(E^{n})$, which is Galois over $K((E/\pm)^{n})=L$ with group $(\ZZ/2)^{n}$:
it is the fibre product of $n$ copies of a degree-$2$ Galois cover, so its deck
group is $(\ZZ/2)^{n}$ by Artin's theorem, with no square-class computation and
in every characteristic. The diagonal fixes every root since $x(R)=x(-R)$. Nothing
else does: a further coincidence among the signed sums forces
$\sum_{i\in S}\epsilon_iP_i\in E[2]$ for some proper nonempty $S$; the summation map
$E^{n}\to E$ is a surjective homomorphism, so the preimage of the finite scheme
$E[2]$ is a proper closed subscheme, and the generic point of the integral variety
$E^{n}$ lies in no proper closed subscheme. The group is abelian, so stabilisers
are normal and the root field is the splitting field. The last step is the whole
content: an earlier sketch of this result omitted it, and a correction records the
gap and its repair (\ev{E12}). The Frobenius element of a good specialisation is the
class of the vector of fibre rationalities modulo global sign, and a non-identity
element of a regular permutation group is fixed-point free, which gives (B). (C)
holds because a rational point makes the decomposition group trivial.

\paragraph{Toy-scale verification.}
The verification run (\ev{E11}, seed $20260807$) found on the factor-base locus
that all $7908$ good specialisations at $m=3$, all $7038$ at $m=4$ and all $1161$
at $m=5$ split completely with roots equal to the signed sums, with no
counterexample ($8000$, $8000$ and $2000$ samples, of which $92$, $962$ and $839$
were degenerate). At $m=3$ it saw $7786$ good generic specialisations, exactly two
Frobenius classes, only the patterns predicted, and maximum factor degree $2$; the
discriminant identity $\mathrm{disc}_T S_3=16f(x_1)f(x_2)$ holds exactly over
$\ZZ[x_1,x_2,A,B]$. Two null objects of the same shape obey the dichotomy in only
$13.69\%$ and $48.91\%$ of separable degree-$4$ samples, so the instrument can
reject (\ev{E11}, \texttt{C7\_null\_controls}).

\paragraph{Status.}
The derivation is complete and self-contained but has not been read by a reader
experienced in function fields; the two points that most need it are Artin's theorem
for a quotient by a finite group acting faithfully, and properness of the preimage
of $E[2]$ in characteristic $2$ where $E[2]$ is non-reduced (\ev{E34}, ``Two blockers''). A separately derived hypothesis reaches the same group, and its toy-scale
execution received a validator and a red-team review (\ev{E13}); that corroborates
the statement and is not a proof review. Novelty has not been adjudicated.

\paragraph{What S4 does not say.}
It is not an algorithm. Statement (C) is about solving for the \emph{last}
coordinate given the other $m-1$; relation search asks the opposite question, with
the target fixed and the other coordinates to be found in the factor base, and
complete splitting gives it no leverage. Its negative content is that a census of
factorisation types or Frobenius classes of the summation fibre is constant on the
factor-base locus and so cannot detect a relation-rate anomaly there, and that no
curve has smaller monodromy. For $m\ge4$ the discriminant is a square in $L$, so a
statistic built on its character is vacuous; $m=3$ is the only $m$ where it carries
information. It implies no speedup, no security reduction and no change to any
complexity estimate.

\section{S5: Lemma 4 of Nagao's ePrint 2015/984 is false as a universal statement}
\label{sec:s5}

\begin{definition}[Definitions 5 and 6 of \cite{Nag15}, equality reading]
\label{def:fall}
Let $f_1,\dots,f_M\in\FF_2[X_1,\dots,X_N]$. A \emph{fall at degree $d$} is a tuple
$(g_i)$ with $\max_i\deg(g_if_i)=d$, $\deg(\sum_ig_if_i)<d$ and $\sum_ig_if_i\neq0$.
The \emph{first fall degree} $d_F$ of a generating set is the least $d$ at which a
fall exists. The \emph{fake} first fall degree $d'_F$ of $\{f_i\}\cup S_{fe}$, with
$S_{fe}=\{X_j^2+X_j\}$, is defined the same way with every degree taken after
reduction modulo $S_{fe}$ and with $\sum_ig_if_i\not\equiv0$.
\end{definition}

Lemma 4 of \cite{Nag15} asserts $d_F\le d'_F$; the paper derives it from a lemma
whose proof it declines to reproduce, and the fake quantity is what its applications
bound (\ev{E29}). The printed condition reads ``$\max_i\deg(g_if_i)\ge d$''; under
that literal reading the lemma holds only vacuously (\ev{E15}). We state the result
for the equality reading, which is the one the complexity estimates need. We
consider $d_F$ with the field equations \emph{outside} the generating set
($F$ alone) and \emph{inside} ($F\cup S_{fe}$).

\begin{proposition}[S5, \textsc{exact}]
\label{prop:lemma4}
Over $\FF_2$ the statement $d_F\le d'_F$ is false in both placements. Witness~A:
$f_1=X^2Y$, $f_2=XY+X$ in $\FF_2[X,Y]$ has $d'_F=2$ and $d_F=3$ outside and inside.
Witness~B: $f_1=X_1X_2+X_3$, $f_2=X_1X_3+X_2$ in $\FF_2[X_1,X_2,X_3]$ has $d'_F=2$ and
$d_F=3$ outside and inside.
\end{proposition}

\begin{proof}[Hand proof for Witness A]
\emph{$d'_F=2$.} Reduced, $f_1\equiv XY$, $f_2=XY+X$. Take $g_1=g_2=1$: both products
have degree $2$ and the sum is $X$, of degree $1<2$ and nonzero, a fall at $2$. In
the reduced ring with basis $1,X,Y,XY$, the multiples $g_1\bar f_1$ lie in $\{0,XY\}$
and the multiples $g_2f_2$ in $\{0,XY+X\}$, since $X\cdot f_2=f_2$ and
$Y\cdot f_2=XY\cdot f_2=0$. A nonzero product has degree $2$, so no tuple has maximum
degree $1$ and no fall occurs below $2$.
\emph{$d_F=3$.} At degree $3$, $g_1=1$, $g_2=X$ gives $f_1+Xf_2=X^2Y+X^2Y+X^2=X^2$,
with both products of degree $3$ and the sum of degree $2$ and nonzero: a fall. For
$d\le2$, $f_1$ has degree $3$ so $g_1=0$; $f_2$ has degree $2$, so at $d=2$ only
$g_2$ constant, and $cf_2$ has degree $2$. Adjoining $X^2+X$ and $Y^2+Y$ changes
nothing: at $d=2$ the leading forms $XY,X^2,Y^2$ of $f_2,X^2+X,Y^2+Y$ are linearly
independent, so no nonzero combination has degree $<2$, and no generator has degree
$\le1$. \qedhere
\end{proof}

\paragraph{Native re-derivation.}
The instrument of \ev{E19} computes $d_F$ and $d'_F$ exactly. For degree $d$ the
tuples with $\deg(g_if_i)\le d$ form an $\FF_2$-vector space $W_d$; those whose
degree-$d$ part of the sum vanishes form a subspace $K$; $Z\subseteq K$ is the
subspace with sum zero. A fall at $d$ exists exactly when $K\not\subseteq Z$ and
$K\not\subseteq W_{d-1}$, because a vector space is never a union of two proper
subspaces. For $d'_F$ the multipliers may be taken multilinear. The output
(\ev{E18}) is in Table~\ref{tab:l4}.

\begin{table}[t]
\centering\small
\begin{tabular}{@{}l c c c@{}}
\toprule
system & $d_F$ outside & $d_F$ inside & $d'_F$ \\
\midrule
control: $X$, $X+1$ (one variable)          & 1 & 1 & 1 \\
Witness A: $X^2Y$, $XY+X$                   & 3 & 3 & 2 \\
Witness B: $X_1X_2+X_3$, $X_1X_3+X_2$       & 3 & 3 & 2 \\
Witness C: $X_1X_2+X_3$, $X_1+X_2$          & 3 & 2 & 2 \\
\bottomrule
\end{tabular}
\caption{First fall degrees from the native instrument (\ev{E18}). Witness C is
the ledger's first, outside-form example: it violates $d_F\le d'_F$ only outside.}
\label{tab:l4}
\end{table}

The instrument also enumerates all $7875$ unordered pairs of distinct nonzero
multilinear polynomials of degree $1$ or $2$ in three variables with a cap of $6$ on
the degree. It finds $d_F>d'_F$ in $6138$ pairs outside and $4830$ inside, and
$d_F<d'_F$ in none outside and $224$ inside; no pair was undetermined at the cap.
The population is an arbitrary finite one and the rates describe it only; they show
that violations are not rare in it. They are not comparable with the earlier
sampled figure of $94$ in $400$ random multilinear pairs (\ev{E14}), which used a
different generator.

\paragraph{Provenance.}
Each witness was computed by two implementations written without sight of each
other (\ev{E14}); the refutation of the \emph{inside} form was first recorded as a
correction of an earlier belief that it held (\ev{E16}, \ev{E17}). The native instrument here is a
third implementation, in Rust, from the definitions.

\paragraph{What S5 does not say.}
Nothing about Proposition 5 of \cite{Nag15}, which may hold by another route, and
nothing about descended summation-polynomial systems at any $n$. The note of \ev{E15} argues that Propositions 2 and 3 of \cite{Nag15} exhibit explicit
certificates and so do not depend on Lemma 4, and that the inference ``fake bound implies
true bound'' is recoverable per instance when every product in a fake-degree certificate
already has unreduced degree at most the fake degree. The justification the paper offers
for its Proposition 5, by contrast, runs through Lemma 4 (\ev{E29}); that justification,
not the proposition, is what S5 affects. The first-fall-degree assumption itself
is a separate matter, on which \cite{HKY15} give evidence against. The result is for
characteristic $2$ and the equality reading.

\section{S6: Assumption 1 of Semaev 2015 off the diagonal}
\label{sec:s6}

Assumption 1 of \cite{Sem15} asserts $d_{F4}\le4$ for the Boolean system equivalent
to the paper's equation~(5) at every $2\le t\le m$ on the diagonal $k=\lceil n/m\rceil$,
where $k$ is the dimension of the subspace $V\subseteq\FF_{2^n}$. The programme's
goal record states that its experimental basis is $2300$ MAGMA systems at $n\le21$,
$m\le6$ plus one cell at $n=40$, $m=t=2$, and that the paper's own Section 4.5.1
reports that the regularity degree ``generally exceeds 4 when $k>\lceil n/m\rceil$
though the first fall degree is still 4'' (\ev{E21}).

The frozen design of this paper (\ev{E27}) listed ``Assumption 1 fails off the diagonal'' as a negative
result. \textbf{We could not establish it.} What the committed records support is
three separate facts.

\paragraph{(i) The source's own remark.}
The off-diagonal behaviour is a remark about the regularity degree in the source
paper, not a programme measurement of $d_{F4}$ (\ev{E21}).

\paragraph{(ii) On the diagonal the assumption is not refuted at reachable sizes.}
With $m=2$ and $V$ a uniformly random $\FF_2$-subspace of dimension exactly $\lceil n/2\rceil$,
the first fall degree is $3.00$, flat, for $n=7,\dots,63$, against a shape-matched
random Boolean null at $4.00$; and a sound refutation predicate for $d_{F4}\le4$
returns \textsc{refuted} at every cell $n=5,\dots,21$, which means the assumption's
bound is consistent with every cell reached (\ev{E20}). The same record notes that an
earlier degree ladder had measured Frobenius-invariant subspaces of dimension
$\mathrm{ord}_n(2)$, which coincide with $\lceil n/2\rceil$ in $0$ of $23$ cells, so
that ladder says nothing for or against the assumption. First fall and solving degree
are not $d_{F4}$, and the record claims no refutation.

\paragraph{(iii) One off-diagonal closure draw reads $D>4$.}
Campaign item 07 measures the degree-4 closure of the chained $S_3$ descent system
with a memory-blocked engine (\ev{E22}). Writing $\delta=l-\lceil n/m\rceil$ for the
subspace dimension $l$ of a cell $(n,l,m)$ relative to the diagonal, its measured
cells are:

\begin{center}\small
\begin{tabular}{@{}lrcp{7.4cm}@{}}
\toprule
cell $(n,l,m)$ & $\delta$ & draws & reading \\
\midrule
$(13,4,4)$ & $0$  & 5 & $D=4$ on all \\
$(11,4,4)$ & $+1$ & 4 (+1 unreceipted) & $D=4$ on all \\
$(13,5,4)$ & $+1$ & 2 & $D=4$ on both \\
$(15,5,4)$ & $+1$ & 3 relayed & one draw saturates without $1$, $D>4$; two stop at the time limit, unmeasured \\
$(19,6,3)$ & $-1$ & 4 & $D=4$ on all \\
\bottomrule
\end{tabular}
\end{center}

Of the three cells with $\delta=+1$, two read $D=4$ on every draw; the third, $(15,5,4)$
with $N=50$ unknowns, has one measured draw of four planned that reads $D>4$. The record itself calls this
a one-draw observation, a closure reading rather than a literal F4 reading (no msolve
reading exists at that cell), and a result at toy scale; the cell is
\textsc{blocked} on time (\ev{E22}).

\begin{proposition}[S6, \textsc{not established}]
\label{prop:s6}
No committed record contains an F4 degree certificate $d_{F4}>4$ at any cell, on or
off the diagonal. The programme's findings state that Assumption 1 is not refuted
(\ev{E14}, ``reads this before using anything below''), and the only off-diagonal
reading above $4$ is one closure draw in item 07. Items 04 and 05 (\ev{E36}, \ev{E37}) measure closure degrees of related descent systems at toy cells, the first with the chained $S_3$ route at $D=4$ on four cells, the second with grevlex at $D_{\mathrm{mat}}=4$ on its countable targets; neither is a reading at the deciding cell.
\end{proposition}

\paragraph{Why it matters, and what would settle it.}
The campaign's pricing of binary curves at deployed sizes (item 02) is conditional on
constant-degree arms ($D=3,4,6,7$) that are hypotheses at deployed $n$, and its gate on
item 07 is open (\ev{E26}); the one arm that is the measured growth fit, $D=l+4$ at
$m=4$, never reaches rho (gap $+105.1$ bits at $n=131$, rising to $+336.7$ at $n=571$,
Koblitz, weight $1$, $K=1$). Settling S6 needs a degree-5 closure at $(15,5,4)$, which
the record puts at $2.37\times10^{6}$ columns and $3.5\times10^{11}$ bytes of reduced
basis, beyond this host, or a different engine. Whether the repository's chained
$S_3$ descent system with a random subspace of dimension $l$ coincides with
Semaev's equation~(5) up to his conventions is the open definitional question
KN-OPEN-d218ec; the sign of $\delta$ above is parameter arithmetic, not a statement
that the two objects agree.

\section{Campaign results that bear on the statements}
\label{sec:campaign}

The blind-spot campaign of 2026-10-18 measured twenty directions. Table~\ref{tab:campaign}
lists the ones whose results bear on S1--S6, quoted from each item's own verdict. Items
03, 06, 08, 09, 11, 12, 13, 15 and 16 concern solver, learned-policy, rho-walk, genus-3 linear-algebra
and toy-size oracle engineering and bear on none of the six; we cite no number from them. The
selection of rows was fixed before the edit (\texttt{PROTOCOL.md}, Addendum B), and each row
quotes the item's own verdict line and level, as the READMEs read when the ledger digests were taken
(item 10's README was still a stage-2 assessment; its frozen rule had no input).

\begin{table}[t]
\centering\small
\begin{tabular}{@{}l p{6.6cm} p{4.3cm}@{}}
\toprule
item & verdict as recorded & bears on \\
\midrule
01 \ev{E24} & \textsc{regresses} at the constant level; \textsc{inconclusive} at the exponent level. No arm cheaper than matched rho at $\log_2\#E=24$ to $40$; best arms $50.0$ [$42.0,60.3$] and $51.6$ [$43.3,62.0$] times matched rho at $40$ bits, weight $0$ & S1 context: distance to rho remains \\
02 \ev{E26} & \textsc{advances} (derived, model level), conditional on a constant-degree model; item-07 gate open; not a measured improvement & S6: the constant-degree hypothesis is the unsettled boundary \\
04 \ev{E36} & \textsc{advances}, level primitive, stage diagnostic: the M\"obius-frame arm lowers the Macaulay closure degree of the incumbent presentation on $5$ of $5$ cells (for example $D$ from $7$ to $5$ at $(17,4,3)$) and does not go below the chained $S_3$ arm, $D=4$ on every cell where it was measured; no $S$, rho ratio or GAE is computed & S6 context: closure degrees at the cells Assumption 1 concerns \\
05 \ev{E37} & \textsc{advances}, level constant, one arm (\texttt{wdeg}): $3.4\times$ to $18.6\times$ fewer matrix-entry touches than this engine's grevlex on four cells; the solving degree is not lowered ($D_{\mathrm{mat}}\ge4$ everywhere) and at the next size the arm hit the time cap & S6 context: grevlex closure degree is $4$ on every countable target \\
07 \ev{E22} & \textsc{blocked} at $(15,5,4)$, level primitive; one draw reads $D>4$ & S6 (iii) \\
10 \ev{E23} & \textsc{inconclusive}, level constant; the frozen rule has no input because the held-out, finalist and conservation cards were not run & S3: no evidence from this item \\
14 \ev{E25} & \textsc{regresses}, level primitive; no $m=4$ oracle's $T_{\mathrm{dec}}$ grows slower than $p^{1/4}$ at the measured scale & S1 context: first measured $m=4$ cells \\
17 \ev{E35} & \textsc{trades}, level constant: the cyclotomic base is never cheaper than the incumbent by the frozen rule at $n=83$ ($S_{\mathrm{ready}}$ ratio $0.995$ to $2.003$ across accountings and cells; about $96\%$ of $S_{\mathrm{ready}}$ is the arm-independent pair table); model rows put the cyclotomic arm at $231.9$ to $1.26\times10^{7}$ times the rho floor at $c=0$ (upper bounds, model rows, not run); the solving-degree leg is \textsc{blocked} & S1 context: distance to rho; S3 context: base-dependent cost is a few percent, yield not measured \\
\bottomrule
\end{tabular}
\caption{Campaign items that bear on the six statements. None is a speedup claim.}
\label{tab:campaign}
\end{table}

\section{Formalisation}
\label{sec:lean}

Two combinatorial cores are stated in Lean~4 with Mathlib, in a new project
\texttt{lean/} (beside the register of \ev{E28}) pinned to the same toolchain (\texttt{v4.34.1}) and Mathlib revision
as the repository's \texttt{S3Roots} project. Table~\ref{tab:proofmap} is the proof map.

\begin{table}[t]
\centering\small
\begin{tabular}{@{}p{4.1cm} p{4.5cm} p{3.9cm}@{}}
\toprule
statement & Lean declaration & status \\
\midrule
Lemma~\ref{lem:core}(i) & \texttt{IcNegative.agree\_card\_le\_one} & written, \textbf{not yet kernel-checked} \\
Lemma~\ref{lem:core}(ii) & \texttt{badSecrets\_card\_le}, \texttt{badSecrets\_card\_le\_sq} & written, not yet kernel-checked \\
Lemma~\ref{lem:core}(iii) & \texttt{injOn\_of\_not\_bad} & written, not yet kernel-checked \\
Prop.~\ref{prop:yield}, any map & \texttt{sum\_yield}, \texttt{sum\_yield\_cast}, \texttt{total\_yield\_eq} & written, not yet kernel-checked \\
coverage bound & \texttt{coverage\_le\_card\_dom}, \texttt{coverage\_le\_card\_cod} & written, not yet kernel-checked \\
Prop.~\ref{prop:yield}, list form & \texttt{FactorBaseDesigns.}\allowbreak\texttt{sum\_count}, \texttt{patterns\_neutral} & kernel-checked in CI \ev{E7}, \ev{E8} \\
Theorem~\ref{thm:151} in full & --- & open (O1) \\
$\mathrm{Sym}^m$ binomial count & --- & open (O2) \\
Prop.~\ref{prop:lemma4} & --- & open (O3); finite check is native \\
Theorem~\ref{thm:mono} & --- & open (O4) \\
\bottomrule
\end{tabular}
\caption{Proof map. ``Not yet kernel-checked'' means the proofs were written on a
host with no Lean toolchain and no Mathlib cache; the repository proof gate is the
first check.}
\label{tab:proofmap}
\end{table}

We do not call the new modules machine-checked. The authoring host had neither a
Lean toolchain nor the Mathlib cache and too little disk to fetch them, so the proofs
were written against the Mathlib interface and have not been elaborated. The
repository's proof gate (\texttt{tools/lean/proof-gate.sh}) builds the project, fails
any theorem whose transitive axioms include \texttt{sorryAx}, and accepts only
\texttt{propext}, \texttt{Classical.choice} and \texttt{Quot.sound}. The sources contain
no \texttt{sorry}, \texttt{admit}, \texttt{native\_decide} or user axiom, and open
targets are kept in \texttt{lean/OPEN\_TARGETS.md} rather than in compiled files.
Until that gate passes, the Lean statements are specifications of the lemmas, whose
mathematical proofs are those given above. The mathematical content is independent of
the formalisation; if a Lean proof needs repair the statement stands.

A Lean file is a statement, and its fidelity to the claim must still be read:
Lemma~\ref{lem:core} is the counting step of Theorem~\ref{thm:151}, not the theorem,
and \texttt{sum\_yield} is the double count behind Proposition~\ref{prop:yield}, not
the multiset instance.

\section{Limitations}
\label{sec:limits}

\begin{itemize}
\item Everything measured is at toy scale; no statement concerns a deployed curve,
  and none reaches the sizes ($n\ge131$) at which a speedup would matter.
\item S1 is a comparison in probes and table entries against bounds stated in
  group operations and bits; it is not a proof of genericity, and the structured
  generic-group bound is vacuous for these rungs.
\item S4 has not had an independent review of its proof, and novelty is not adjudicated.
\item S5 is for characteristic $2$ and the equality reading and says nothing about
  Proposition 5 of \cite{Nag15}.
\item S6 is open. We include it so that a reader does not count it among the
  closed negatives.
\item Several evidence files were produced by sibling repositories whose heads moved
  during the work; Appendix~\ref{app:ledger} fixes the digests read.
\item The new Lean modules have not been kernel-checked (Section~\ref{sec:lean}).
\end{itemize}

\appendix
\section{Evidence ledger}
\label{app:ledger}

Each row is a file read for this paper, with the first $16$ hex digits of its SHA-256.
The full digests are in \texttt{evidence/ledger.full}; the ledger is regenerated by
\texttt{evidence/make\_ledger.sh}. Repository heads when generated are in
\texttt{EVIDENCE.md}.

{\footnotesize
\begin{longtable}{@{}l p{10.6cm} l@{}}
\toprule
id & repository:path & sha256 (16) \\
\midrule
\endhead
E1 & \nolinkurl{crypto:research/lit2026_w3_ggm_audit_20261011/reports/run-20261011T2012Z/table.json} & \texttt{26d462fe6f6de45c} \\
E2 & \nolinkurl{crypto:research/lit2026_w3_ggm_audit_20261011/reports/run-20261011T2012Z/report.md} & \texttt{2062646518f95046} \\
E3 & \nolinkurl{crypto:research/lit2026_w3_ggm_audit_20261011/README.md} & \texttt{4a023d24168ddda6} \\
E4 & \nolinkurl{crypto:RESEARCH_ISOGENY_FIBERS_ECDLP.md} & \texttt{23205b00c7cc4c25} \\
E5 & \nolinkurl{crypto-autoresearcher:knowledge/findings/KN-FIND-007.md} & \texttt{9dd306d2fad7387d} \\
E6 & \nolinkurl{crypto:research/cyclotomic_falsifier_20261015/README.md} & \texttt{be5199e1a20534e0} \\
E7 & \nolinkurl{crypto:research/factor_base_designs_20261015/lean/FactorBaseDesigns/Yield.lean} & \texttt{e325cb42e2c560fc} \\
E8 & \nolinkurl{crypto:research/factor_base_designs_20261015/proof_logs/axioms.txt} & \texttt{06148bb5c423bf9f} \\
E9 & \nolinkurl{crypto-autoresearcher:knowledge/findings/KN-FIND-a8990a.md} & \texttt{aaefb551e0fbee1c} \\
E10 & \nolinkurl{crypto-autoresearcher:papers/semaev-conservation-specialization/paper.tex} & \texttt{6ce58081151c6e06} \\
E11 & \nolinkurl{crypto-autoresearcher:experiments/EXP-SMON-e5cbe6/runs/RUN-SMON-e5cbe6-001/raw-result.json} & \texttt{a7f3cb1265cf1102} \\
E12 & \nolinkurl{crypto-autoresearcher:ledger/corrections/CORR-20260807-652652.yaml} & \texttt{67da07f933b70565} \\
E13 & \nolinkurl{crypto-autoresearcher:knowledge/findings/KN-FIND-6edecd.md} & \texttt{935ede09e62a5abf} \\
E14 & \nolinkurl{crypto-autoresearcher:knowledge/findings/KN-FIND-936151.md} & \texttt{a17623ab3189cba0} \\
E15 & \nolinkurl{crypto-autoresearcher:papers/nagao-lemma4-note/note.md} & \texttt{300372bab1ebd7fc} \\
E16 & \nolinkurl{crypto-autoresearcher:ledger/corrections/CORR-20260921-942a62.yaml} & \texttt{f9def986c5ec51cc} \\
E17 & \nolinkurl{crypto-autoresearcher:tools/lemma4_inside_form.py} & \texttt{779e865265a768a7} \\
E18 & \nolinkurl{crypto:papers/ic-negative-results/evidence/checks.txt} & \texttt{ddd155bb68e11e12} \\
E19 & \nolinkurl{crypto:examples/blindspots_19_checks.rs} & \texttt{751277a8c075483c} \\
E20 & \nolinkurl{crypto-autoresearcher:knowledge/findings/KN-FIND-9cba50.md} & \texttt{812688a01a31d18a} \\
E21 & \nolinkurl{crypto-autoresearcher:ledger/goals/GOAL-SEMBIN-5078bc/goal.yaml} & \texttt{107a5cc6507e58ff} \\
E22 & \nolinkurl{crypto:research/blindspots_20261018/07-blocked-f4-degree-cells/README.md} & \texttt{131ad266f6cb604b} \\
E23 & \nolinkurl{crypto:research/blindspots_20261018/10-kernel-program-search/README.md} & \texttt{67dd70829c4e08a7} \\
E24 & \nolinkurl{crypto:research/blindspots_20261018/01-fp2-gaudry-k2/README.md} & \texttt{34d06d7bdec87acc} \\
E25 & \nolinkurl{crypto:research/blindspots_20261018/14-prime-m4-cells/README.md} & \texttt{5a333ca46a54939f} \\
E26 & \nolinkurl{crypto:research/blindspots_20261018/02-deployed-binary-pricing/README.md} & \texttt{529c38f03889ab2b} \\
E27 & \nolinkurl{crypto:research/blindspots_20261018/EXPERIMENTS.md} & \texttt{fe91cf9a1abd32fb} \\
E28 & \nolinkurl{crypto:research/factor_base_designs_20261015/README.md} & \texttt{7aec9d11222cf315} \\
E29 & \nolinkurl{crypto-autoresearcher:coordination/review/sembin-20260916-propquant/report.md} & \texttt{a0c7642720eef546} \\
E30 & \nolinkurl{crypto:research/blindspots_20261018/11-learned-prefilter/README.md} & \texttt{943422dff35ed15f} \\
E31 & \nolinkurl{crypto:papers/ic-negative-results/lean/IcNegative/AffineCore.lean} & \texttt{178ea2dd12090590} \\
E32 & \nolinkurl{crypto:papers/ic-negative-results/lean/IcNegative/YieldConservation.lean} & \texttt{ee973d4cd005356f} \\
E33 & \nolinkurl{crypto:research/pair46_factor_bases_20261015/13-xy-windows/RESULTS.md} & \texttt{f2f6a266dc807687} \\
E34 & \nolinkurl{crypto-autoresearcher:papers/semaev-conservation-specialization/README.md} & \texttt{99888c7d58c128e6} \\
E35 & \nolinkurl{crypto:research/blindspots_20261018/17-cyclotomic-cost-n83/README.md} & \texttt{7ba9a2653dce7a3d} \\
E36 & \nolinkurl{crypto:research/blindspots_20261018/04-presentation-search/README.md} & \texttt{d4676ae4b5ec3e4b} \\
E37 & \nolinkurl{crypto:research/blindspots_20261018/05-elimination-orders/README.md} & \texttt{845129e7617bb62d} \\
\bottomrule
\end{longtable}

}

\section{Reproducing the native checks}
\label{app:repro}

\begin{verbatim}
rustc --edition 2021 -C opt-level=3 -o b19 examples/blindspots_19_checks.rs
prlimit --as=4294967296 ./b19 all > output.txt   # about 0.3 s, one thread
diff <(tail -n +3 output.txt) \
     <(tail -n +3 papers/ic-negative-results/evidence/checks.txt)
\end{verbatim}
\noindent The output ends with \texttt{ALL CHECKS PASS}. The subcommands are
\texttt{lemma4}, \texttt{affine} and \texttt{yield}.

\begin{thebibliography}{99}
\bibitem{BL13} D.~J. Bernstein and T.~Lange. Non-uniform cracks in the concrete: the
power of free precomputation. ASIACRYPT 2013.
\bibitem{CGK18} H.~Corrigan-Gibbs and D.~Kogan. The discrete-logarithm problem with
preprocessing. EUROCRYPT 2018.
\bibitem{DKM25} I.~Dinur, N.~Keller and O.~Marmor. arXiv:2505.00894 (STOC 2026),
Theorem 1.1, as quoted in \ev{E3}.
\bibitem{HKY15} M.-D.~A. Huang, M.~Kosters and S.~L. Yeo. Last fall degree, HFE, and
Weil descent attacks on ECDLP. CRYPTO 2015 (ePrint 2015/573).
\bibitem{Nag15} K.~Nagao. Complexity of ECDLP under the first fall degree assumption.
ePrint 2015/984.
\bibitem{Sem15} I.~Semaev. New algorithm for the discrete logarithm problem on
elliptic curves. ePrint 2015/310.
\bibitem{Sho97} V.~Shoup. Lower bounds for discrete logarithms and related problems.
EUROCRYPT 1997.
\end{thebibliography}

\end{document}
